\section*{Conclusions}

These document has covered the development and evolution of a Heat and Mass Transfer solver across different benchmarks of increasing complexity. Across the different stages of development of this model, different complexities regarding the implementation of numerical methods and the simplification of physical models have been observed and analyzed. In addition, the model developed during this project will serve as a stepping stone towards more complex implementations. There are still several improvements that can be made to the model, particularly in terms of optimization of variables and standarization of certain algorithms. The model will continue to be maintained to expand its capacity and increase its performance.

\begin{enumerate}
    \item The Pure Diffusion Model served as a first step in defining the architecture of the model. This structure consists of the 5 main components (Material, Mesh, Discretizer, Solver, and Probe) that handle the different operations required for the solver; storing data about the domain, about the physical phenomena, solving the model, and storing the desired results. This system allows for organized and easy-to-read code, which makes it easier to work on long-term maintenance of the model. Furthermore, this model was used as a baseline comparison for the linear solvers implemented in the model (GS and CG).
    \item The Convection Diffusion Model served as an introduction to the advective scheme, which determines the numerical approach for approximating the convective effect of flow at the walls. The formulation proposed by Patankar allows for a standardized model that can handle different advective schemes under a single formulation. It is important to highlight the observed behaviour of CDS when used to solve convection-dominant problems, as this model can lose diagonal dominance for $|Pe| > 2$, showing a checkerboarding effect. For high $Pe$ cases on the other hand, UDS can provide a more stable solution while sacrificing precision due to gradient smoothing. 
    \item The Navier-Stokes Model introduced the Fractional Step Method, which consists of separating the velocity components from the pressure component and solving them discretely by using the Adams-Bashforth discretization method to obtain second order accuracy. Furthermore, Staggered Meshes were used to eliminate the need for interpolating values at the faces; choosing to locate them at the required locations from the start. Finally, the difference between a diffusion-dominant and convection-dominant scenario was also evident in this tests. However, since the velocity is decoupled and calculated explicitly with Adams-Bashforth, the same checkerboarding apparent for the convection-diffusion case was not seen here.
    \item All models were succesfully validated against the benchmark reference values. 
\end{enumerate}


\subsection{Future Work}

\begin{enumerate}
    \item During the development of the model, the code architecture has changed significantly. As the model kept expanding the, new variables and structs were introduced to store the required information and better organize the code. While there were several changes in structure when moving from a pure-diffusion model to a convection-diffusion model, the main changes (particularly to the Mesh) were made when moving to the Navier-Stokes model. Since the first two models solved a single transport variable, Mesh.cpp stored all relevant variables directly. On the other hand, the NS solver requires multiple transport variables to be evaluated at the same time; and as such, they require their own set of variables for that purpose. For this purpose, a cleaner struct implementation was developed that organizes each transport variable independently and stores different metadata to reduce the operations done by the discretizer. The convection-diffusion solver should be modified to operate under the same parameters, modifying the logic to fit the new struct definition and standardizing the model.
    \item The current model can still be expanded to support more complex cases, with the next step being solving natural convection problems. This requires the inclusion of the energy conservation equation in the model as well. In addition, the current model can still be extended to support a 3D domain; however, this change implies a larger overhaul of several functions in the model and will be considered for implementation at a later date.
    \item Further studies about the nature of the numerical solver can be included to deepen the understanding and implication of this decisions. For instance, only CDS and UDS were tested during the development of this project even thought the Hybrid, Exponential, and Power Law schemes are also implemented. With more time and computational resources, additional tests for the 5 different schemes could be done in different edge scenarios.
\end{enumerate}
